% TOP_PROBLEM: {"id": 6700064, "title": "Scalar Curvature Question [?68]: Let $\\widetilde X$ be the universal cover of a Riemannian $n$-manifold $X$ homeomorphic to the $n$-torus", "category_id": 6, "category": {"id": 6, "name": "geometry", "display_name": "Geometry", "description": "Euclidean and non-Euclidean geometry, geometric structures.", "slug": "geometry", "order_index": 6, "created_at": "2026-07-31T15:26:25.671Z"}, "status": "partially_solved"}
% FIRST_POSTED: null
% ENTRY_KIND: statement_only
% Imported from: batch16.tex; UnsolvedMath ID 6700064
\documentclass[11pt]{article}
\usepackage[margin=1in]{geometry}
\usepackage{amsmath,amssymb,mathtools}
\usepackage{fontspec,unicode-math}
\setmainfont{Latin Modern Roman}
\setsansfont{Latin Modern Sans}
\setmathfont{Latin Modern Math}
\usepackage{xurl,hyperref}
\hypersetup{colorlinks=true,linkcolor=blue,urlcolor=blue}
\newcommand{\sref}[2]{\hyperlink{source:#1:#2}{[S#2]}}
\newcommand{\cataloguescope}[1]{\paragraph{Mathematical question.}#1}
\begin{document}
% UnsolvedMath ID 6700064; source catalogue code AMR-066-0064
\setcounter{section}{6700063}
\section{Scalar Curvature Question [?68]: Let \$\textbackslash{}widetilde X\$ be the universal cover of a Riemannian \$n\$-manifold \$X\$ homeomorphic to the \$n\$-torus}\label{6700064}
{\small\sffamily UnsolvedMath ID 6700064 · Geometry}
\subsection{Definitions and mathematical statement}

\paragraph{Shared notation.}$\mathbb N=\{1,2,\ldots\}$, and $\mathbb Z,\mathbb Q,\mathbb R,\mathbb C$ denote the integers, rationals, reals and complex numbers. Any other conventions are specified locally.


For a Riemannian domain, mean curvature $H$ of a smooth boundary face is the trace of its second fundamental form with the convention that the boundary of a Euclidean ball has positive mean curvature. Mean convexity means $H\geq0$. Dihedral angles are interior angles between adjacent faces, measured between their tangent spaces along their common codimension-two faces. Let $(X,g)$ be a Riemannian $n$-manifold homeomorphic to the torus $T^n=(S^1)^n$, and let $\widetilde X$ be its universal cover with the lifted metric. An exhaustion by domains is a sequence of compact domains $P_i\subset\widetilde X$ whose interiors exhaust $\widetilde X$ and with $P_i\subset P_{i+1}$. An overcubic domain has a corner structure and a degree-one map $f_i:P_i\to[0,1]^n$ respecting faces in the corner-proper sense: for every codimension-one cube face $\bar F_j$, its corresponding domain face is exactly $F_j=f_i^{-1}(\bar F_j)$. Consequently each lower-dimensional face is sent to its corresponding cube face. Degree is the integer by which the induced map multiplies the oriented relative fundamental class in top-dimensional homology. The phrase ``non-positive scalar curvature at infinity'' in this question refers to the proposed exhaustion below, without a pointwise sign assumption on the scalar curvature of $g$.
\paragraph{Statement recorded in the supplied source.}
Prove that $\widetilde X$ admits an exhaustion by overcubic domains $P_i$ with corners such that every codimension-one face has strictly positive mean curvature and every dihedral angle is at most $\pi/2$. \sref{6700064}{2}
\subsection{Short English statement}
Can the universal cover of any Riemannian torus be exhausted by domains mapping with degree one to a cube, with positive face mean curvature and no dihedral angle larger than a right angle?
\subsection{Sources}
\begin{description}
\item[\hypertarget{source:6700064:1}{S1}] ULAM.AI and the UnsolvedMath Contributors, UnsolvedMath: A Curated Collection of Open Mathematics Problems, record 6700064 (AMR-066-0064). CC BY 4.0. \url{https://huggingface.co/datasets/ulamai/UnsolvedMath}.
\item[\hypertarget{source:6700064:2}{S2}] Misha Gromov, \emph{Four Lectures on Scalar Curvature}, arXiv:1908.10612v6 (9 July 2021), \S3.13, Exercise and full hint, printed p. 168; corner-map definitions in \S3.18, printed p. 186. \url{https://arxiv.org/abs/1908.10612v6}.
\end{description}
\label{end:6700064}
\end{document}
